\section{制度需要真实 Steering Wheel}

访谈最后把政府与企业的失灵解释为 feedback 断裂：领导者拥有 strategy 和 dashboard，却无法及时看到执行现场；一线人员认为上层无知，又把真实工作搬到 Excel；几个月后才出现聚合报告，错误在层级之间累积。Sankar 用 Disneyland Jungle Cruise 的假方向盘比喻这种状态。软件若想恢复机构能力，必须把 intent、execution、outcome 和 learning 接成实时且可审计的 loop。

\begin{figure}[H]
\centering
\includegraphics[width=0.94\textwidth]{images/13-institutional-feedback.png}
\caption{真实 steering wheel 需要目标、执行、结果与学习构成闭环。}
\figsource{依据字幕 42:00--44:03 重绘，`images/13-institutional-feedback.png`。}
\end{figure}

\begin{importantbox}{读图：闭环成立的四个验收问题}
目标是否被转成可执行约束；执行状态是否在同一系统中可见；outcome 是否能关联到具体 decision；学习是否真正修改下一轮 policy 或 workflow。只要其中一步靠离线 spreadsheet、口头会议或不可追溯的人工搬运，方向盘就会逐渐与现场脱离。
\end{importantbox}

\begin{table}[H]
\centering
\begin{tabular}{L{0.19\textwidth}L{0.34\textwidth}L{0.37\textwidth}}
\toprule
断点 & 表面症状 & 应建设的能力 \\
\midrule
Intent→执行 & 员工不知道优先级为何变化 & 版本化 policy、constraint 与 owner \\
执行→结果 & Dashboard 只有活动量，没有结果 & Entity-level lineage 与 outcome metric \\
结果→学习 & 复盘无法定位哪项决策造成后果 & Decision log、counterfactual 和 audit \\
学习→Intent & 同类事故重复发生 & Rule update、workflow test 与 rollout \\
正式系统→Excel & 关键决策离开平台 & 更快建模、可解释接口和一线共创 \\
\bottomrule
\end{tabular}
\caption{制度反馈断点应转成可观察的软件与组织能力。}
\end{table}

\begin{knowledgebox}{课堂提示：不要把“机构失灵”直接推成“推倒重来”}
讲者反对因为现有组织低效就放弃机构本身。系统工程的任务是恢复 sensing、decision 和 feedback，让 policy 能从结果中学习。这个判断同样需要约束：修复机构不能成为扩大不透明权力的理由，改造过程应保留监督、分权和纠错。
\end{knowledgebox}

\subsection{本章小结}

Decision infrastructure 的最终价值不是漂亮 dashboard，而是让组织拥有真实 steering wheel。Intent、execution、outcome 和 learning 必须共享 entity、权限和审计链。Excel escape、延迟报告和无法解释的聚合指标，都是闭环失真的信号。

